\documentclass[11pt]{article}

\usepackage[T1]{fontenc}
\usepackage[utf8]{inputenc}
\usepackage{amsmath,amssymb,amsthm,mathtools}
\usepackage[margin=1in]{geometry}
\usepackage{microtype}
\usepackage[hidelinks]{hyperref}
\emergencystretch=1em
\hypersetup{
  pdftitle={The First Stabilisation Kernel for the Horvath--Kania Motakis Space},
  pdfauthor={Xiaojian Zeng},
  pdfkeywords={Banach operator algebra, Calkin algebra, invertible group, stabilisation}
}

\newtheorem{theorem}{Theorem}
\newtheorem{lemma}[theorem]{Lemma}
\theoremstyle{remark}
\newtheorem{remark}[theorem]{Remark}

\newcommand{\B}{\mathcal B}
\newcommand{\K}{\mathcal K}
\newcommand{\Q}{\mathcal Q}
\newcommand{\Inv}{\operatorname{Inv}}
\newcommand{\Invzero}{\operatorname{Inv}_0}
\newcommand{\GL}{\operatorname{GL}}
\newcommand{\Id}{\mathfrak I}

\title{The First Stabilisation Kernel for the\\Horv\'ath--Kania Motakis Space}
\author{Xiaojian Zeng\\
PBC School of Finance, Tsinghua University, 100190\\
\texttt{xiaojianzeng@foxmail.com}}
\date{}

\begin{document}

\maketitle

\begin{abstract}
Let $X$ be the particular complex Motakis realisation over $S^4$ used by
Horv\'ath and Kania, let $E=X\oplus X$, and let $S,T\in\B(E)$ be their fixed
operators.  We compute the two component groups relevant to their first
stabilisation question:
\[
  \Id(\B(E))\cong \mathbb Z/2,
  \qquad
  \Id(M_2(\B(E)))=1.
\]
The nonidentity element in the first group is $[I_E-2ST]$.  Consequently the
kernel of the first stabilisation is
\(
\{[I_E],[I_E-2ST]\}\cong\mathbb Z/2
\), and $[I_E-2ST]$ has order two.  The main step is a norm-controlled lemma:
for a complex Banach space with a Schauder basis, the Calkin quotient induces
an injection on components of invertibles.  This answers Question~9.1 of the
cited source for exactly the realisation and lifts occurring there.
\end{abstract}

\noindent\textbf{Keywords.}
Banach operator algebra; Calkin algebra; invertible group; stabilisation;
Schauder basis; component group.

\medskip
\noindent\textbf{2020 Mathematics Subject Classification.}
Primary 46H05, 47L10; Secondary 46B25, 55Q15.

\section{The question and the answer}

For a complex Banach space $Y$, write $\B(Y)$ and $\K(Y)$ for its bounded and
compact operators, and put
\[
  \Q(Y)=\B(Y)/\K(Y).
\]
For a complex unital Banach algebra $A$, let $\Invzero(A)$ be the norm-connected
component of $1_A$ in $\Inv(A)$ and set
\[
  \Id(A)=\Inv(A)/\Invzero(A).
\]
The identity component is a normal subgroup, so $\Id(A)$ is the group of
components of the invertible group.  The first stabilisation homomorphism is
\[
 s_A:\Id(A)\longrightarrow\Id(M_2(A)),
 \qquad
 s_A([a])=
 \left[\begin{pmatrix}a&0\\0&1_A\end{pmatrix}\right].
\]
This is an unstable component invariant, rather than the direct-limit group
$K_1(A)$.  The stable $K$-theory of algebras of operators on Banach spaces has
also been studied; see, for example, Laustsen~\cite{Laustsen1997}.

Motakis constructed, for every compact metrisable $K$, a Banach space whose
Calkin algebra is isomorphic to $C(K)$~\cite{Motakis2024}.  Horv\'ath and
Kania use a prescribed-family form of that construction over $S^4$ to obtain a
particular complex space $X$, put $E=X\oplus X$, and construct operators
$S,T\in\B(E)$~\cite[Theorems~4.1 and~5.1]{HorvathKania2026}.  The choices of
realisation and lifts matter throughout this note: $X,E,S,T$ always denote
exactly those objects, and no assertion is made about arbitrary Motakis
realisations or arbitrary lifts.

Their construction lifts the topological obstruction used by Klaja and
Ransford~\cite{KlajaRansford2017}.  That obstruction lies in the invertible
group of $C(S^4,M_2(\mathbb C))$.  If
\[
  u=I_E-2ST,
\]
then $u$ is invertible, its class is not in the identity component of
$\Inv(\B(E))$, and its Calkin image represents the nonzero element of
$\pi_4(\GL_2(\mathbb C))\cong\mathbb Z/2$.  The source also proves that
$[u]$ is killed by first stabilisation, but observes that the order of the
Calkin image does not by itself determine the order upstairs.  Its
Question~9.1 asks for the full kernel and for the order of $[u]$
\cite[pp.~23 and~26]{HorvathKania2026}.

The problem here is topological: it concerns components of invertibles, not a
quantitative functional-calculus estimate.  For comparison with that distinct
operator-theoretic direction, Zeng~\cite{Zeng2026Lens} studies spectral-set
constants for sectorial numerical-range models and determines exact restricted
operator classes.  None of those quantitative results is used below.

Our answer is as follows.

\begin{theorem}\label{thm:main}
For the particular $X,E,S,T$ just specified,
\[
  \Id(\B(E))=\{[I_E],[I_E-2ST]\}\cong\mathbb Z/2
  \quad\text{and}\quad
  \Id(M_2(\B(E)))=1.
\]
Therefore
\[
 \ker s_{\B(E)}
   =\{[I_E],[I_E-2ST]\}\cong\mathbb Z/2,
 \qquad
 \operatorname{ord}[I_E-2ST]=2.
\]
\end{theorem}

The only missing passage from the Calkin calculation to the upstairs
calculation is supplied by the next lemma.  Its proof includes the relevant
norm controls because an arbitrary quotient path need not possess a global
lift.

\section{Injectivity of Calkin components}

\begin{lemma}[Calkin-component injection]\label{lem:injection}
Let $Y$ be an infinite-dimensional complex Banach space with a Schauder basis,
and let
\[
 q:\B(Y)\longrightarrow\Q(Y)
\]
be the quotient map.  Then the induced homomorphism
\[
 q_*:\Id(\B(Y))\longrightarrow\Id(\Q(Y))
\]
is injective.
\end{lemma}

\begin{proof}
Let $(P_n)$ be the partial-sum projections associated with the basis.  They
converge strongly to $I_Y$, and the uniform boundedness principle gives
\[
  \beta:=\sup_n\lVert P_n\rVert<\infty.
\]
We first record explicitly that every compact operator on $Y$ is a norm limit
of finite-rank operators.  Given $K\in\K(Y)$, the set
\(
C=\overline{K(B_Y)}
\)
is compact.  For $\varepsilon>0$, choose a finite
$\varepsilon/(3(1+\beta))$-net $z_1,\dots,z_r$ in $C$.  Strong convergence
allows one to choose $n$ so that
\[
 \max_{1\le j\le r}\lVert(I_Y-P_n)z_j\rVert<\varepsilon/3.
\]
For any $z\in C$, take $z_j$ in the chosen net.  Then
\[
 \lVert(I_Y-P_n)z\rVert
 \le (1+\beta)\lVert z-z_j\rVert
       +\lVert(I_Y-P_n)z_j\rVert
 <\frac{2\varepsilon}{3}<\varepsilon.
\]
It follows that
\[
 \lVert K-P_nK\rVert<\varepsilon,                       \tag{2.1}
\]
and $P_nK$ has finite rank.

We next show that the invertible fibre over $1_{\Q(Y)}$ is path connected.
Take
\(
a=I_Y+K\in\Inv(\B(Y))
\)
with $K$ compact.  By (2.1), choose a finite-rank $F$ satisfying
\[
 \lVert a^{-1}\rVert\,\lVert F-K\rVert<1.              \tag{2.2}
\]
For $0\le t\le1$,
\[
 a+t(F-K)=a\bigl(I_Y+t a^{-1}(F-K)\bigr)
\]
is invertible by the Neumann series, so $I_Y+K$ is joined to $I_Y+F$.
Put $M=F(Y)$ and choose a closed complement $N$, so $Y=M\oplus N$.  Relative
to this decomposition,
\[
 I_Y+F=
 \begin{pmatrix}C&B\\0&I_N\end{pmatrix}
 \quad\text{with}\quad C\in\GL(M).
\]
Indeed, invertibility of the triangular operator forces invertibility of its
finite-dimensional diagonal block $C$.  The path
\[
 t\longmapsto
 \begin{pmatrix}C&(1-t)B\\0&I_N\end{pmatrix}
\]
removes the off-diagonal block through invertibles.  Finally,
$\GL(M)$ is path connected over the complex scalars.  For example, after
choosing a Hermitian structure on $M$, polar decomposition joins $C$ to a
unitary matrix through positive invertible matrices, and a unitary logarithm
joins that unitary matrix to $I_M$.  Extending this path by $I_N$ gives a path
to $I_Y$.  We have therefore proved that
\[
 \{I_Y+K:K\in\K(Y),\ I_Y+K\in\Inv(\B(Y))\}             \tag{2.3}
\]
is path connected.

Now suppose that $v\in\Inv(\B(Y))$ and that $q(v)$ belongs to
$\Invzero(\Q(Y))$.  Choose a continuous path
\[
 \gamma:[0,1]\longrightarrow\Inv(\Q(Y)),
 \qquad \gamma(0)=q(v),\quad \gamma(1)=1_{\Q(Y)}.
\]
The quantity
\(
L=\max_{0\le t\le1}\lVert\gamma(t)^{-1}\rVert
\)
is finite.  By uniform continuity, there is a partition
$0=t_0<\cdots<t_m=1$ such that
\[
 \lVert\gamma(t_{i+1})-\gamma(t_i)\rVert<\frac{1}{2L}
 \quad(0\le i<m).
\]
Thus, for
\(
\delta_i=\gamma(t_i)^{-1}\gamma(t_{i+1})
\),
\[
 \lVert\delta_i-1_{\Q(Y)}\rVert<\frac12.
\]
The definition of the quotient norm supplies $h_i\in\B(Y)$ such that
\[
 q(h_i)=\delta_i-1_{\Q(Y)},
 \qquad \lVert h_i\rVert<1.                            \tag{2.4}
\]
Starting with $v_0=v$, define $v_{i+1}=v_i(I_Y+h_i)$.  The segment
\[
 s\longmapsto v_i(I_Y+s h_i),\qquad 0\le s\le1,
\]
consists of invertibles by (2.4).  Inductively,
$q(v_i)=\gamma(t_i)$; in particular, $q(v_m)=1_{\Q(Y)}$.  Hence
$v_m=I_Y+K_0$ for a compact $K_0$, and (2.3) joins $v_m$ to $I_Y$.
Concatenating the finitely many paths proves
$v\in\Invzero(\B(Y))$.  Thus the kernel of $q_*$ is trivial.
\end{proof}

\begin{remark}\label{rem:hypotheses}
The proof uses the Schauder basis only to obtain the explicit finite-rank
approximation (2.1), and it uses complex scalars in the connectedness of the
finite-dimensional general linear group.  It neither lifts the whole quotient
path nor contracts $I_Y+K$ along the generally unsafe segment $I_Y+tK$.
\end{remark}

\section{Application to the fixed construction}

We now use only properties of the fixed Horv\'ath--Kania construction.  Their
Theorem~4.1 gives $X$ a Schauder basis and a unital Banach-algebra isomorphism
\[
  \Q(X)\cong C(S^4).
\]
Their block-matrix identification gives, for each positive integer $n$,
\[
  \Q(X^n)\cong M_n(C(S^4))
             \cong C(S^4,M_n(\mathbb C)).              \tag{3.1}
\]
These statements, including the prescribed diagonal lifts needed to define
$S$ and $T$, are part of the cited construction
\cite[Theorem~4.1 and Lemma~4.5]{HorvathKania2026}.

Put $u=I_E-2ST$.  Theorem~5.1 of the source proves that $u$ is invertible and
identifies $q(u)$ with the Klaja--Ransford matrix function.  The component
calculation in the same source gives
\[
 \Id(\Q(E))
 \cong \Id(C(S^4,M_2(\mathbb C)))
 \cong \pi_4(\GL_2(\mathbb C))
 \cong \mathbb Z/2,                                    \tag{3.2}
\]
and $q(u)$ represents the nonzero class
\cite[Theorem~5.1 and p.~23]{HorvathKania2026}.  Since $E=X^2$ has a
Schauder basis, Lemma~\ref{lem:injection} makes
\[
 q_*:\Id(\B(E))\longrightarrow\Id(\Q(E))
\]
injective.  Its image contains the identity and, by the existence of the
invertible lift $u$, the nonzero element of (3.2).  It is therefore an
isomorphism, and
\[
 \Id(\B(E))=\{[I_E],[u]\}\cong\mathbb Z/2.             \tag{3.3}
\]

For the stabilised group, the canonical block-operator isomorphism and (3.1)
give
\[
 M_2(\B(E))\cong\B(E\oplus E),
 \qquad
 \Q(E\oplus E)\cong C(S^4,M_4(\mathbb C)).             \tag{3.4}
\]
Polar decomposition retracts $\GL_n(\mathbb C)$ onto $U(n)$.  The stable-range
maps and Bott periodicity~\cite{Bott1957} yield
\[
 \pi_4(\GL_4(\mathbb C))=0;                            \tag{3.5}
\]
this is also the stable-range computation recorded in the source
\cite[p.~2]{HorvathKania2026}.  For completeness, (3.5) indeed controls free
maps from $S^4$: given $f:S^4\to\GL_4(\mathbb C)$ and a base point $x_0$,
the map $x\mapsto f(x_0)^{-1}f(x)$ is based and hence null-homotopic, while
$\GL_4(\mathbb C)$ is path connected.  Thus every such $f$ is homotopic to the
constant identity map.  Pointwise invertibility now gives
\[
 \Id(C(S^4,M_4(\mathbb C)))=1.                         \tag{3.6}
\]

The space $E\oplus E=X^4$ again has a Schauder basis.  Lemma
\ref{lem:injection}, (3.4), and (3.6) therefore imply
\[
 \Id(M_2(\B(E)))\cong\Id(\B(E\oplus E))=1.             \tag{3.7}
\]
Consequently the first stabilisation map has all of (3.3) as its kernel.
Since $[u]$ is the nonidentity element of a two-element group, its order is
exactly two.  This proves Theorem~\ref{thm:main}.

\section*{Scope and assistance disclosure}

The conclusion resolves the two questions posed in Question~9.1 of
Horv\'ath--Kania~\cite{HorvathKania2026} for the exact $X,E,S,T$ constructed
there.  It does not assert a corresponding computation for other Motakis
realisations, real Banach spaces, arbitrary operator lifts, or arbitrary Calkin
extensions.  The general Lemma~\ref{lem:injection} is used only for $X^2$ and
$X^4$ in the application.

OpenAI Codex assisted with source inspection, bibliographic record checking,
proof stress-testing, manuscript drafting, and the deterministic LaTeX build.
No numerical experiment or computer algebra calculation is used as evidence
for the theorem; all mathematical steps needed for the conclusion appear in
the text.

\bibliographystyle{plain}
\bibliography{references}

\end{document}
